\documentclass[11pt]{article}
\usepackage[margin=1in]{geometry}
\usepackage[T1]{fontenc}
\usepackage{lmodern,microtype}
\usepackage{amsmath,amssymb,booktabs,array}
\usepackage[colorlinks,linkcolor=black]{hyperref}
\newcommand{\K}{\mathbb{K}}
\newcommand{\E}{\mathbb{E}}
\newcommand{\F}{\mathbb{F}}
\newcommand{\cX}{\mathsf{X}}
\newcommand{\cY}{\mathsf{Y}}
\newcommand{\cC}{\mathsf{C}}
\newcommand{\cV}{\mathsf{V}}
\newcommand{\bit}{\mathrm{bit}}
\newcommand{\car}{\mathrm{carry}}
\setlength{\parskip}{0.4em}
\newcommand{\head}[1]{\par\vspace{0.35em}\textbf{#1}\enspace}
\raggedbottom
\setlength{\parindent}{0pt}
\title{\vspace{-2.5em}The quotient PIOP}
\author{}
\date{\vspace{-3.5em}}
\begin{document}
\maketitle

A polynomial IOP for circuits made of one gate, a full adder, over a field of characteristic 2. The prover sends two batches of polynomials and five field elements. The verifier sends three random messages, then asks for four evaluations. Nothing here depends on how polynomials are committed.

\section{Setting}

\head{Two fields.} $\K = \mathrm{GF}(2^{64})$ holds the witness. $\E = \mathrm{GF}(2^{192})$, of degree 3 over $\K$, holds the verifier's challenges.

\head{The domain.} Let $S(X) = X^2 + X$ and let $w_k$ be $S$ composed $k$ times:
\[
  w_0(X) = X, \qquad w_{k+1} = w_k^2 + w_k .
\]
Fix $\beta_0 = 1, \beta_1, \dots, \beta_{m-1}$ in $\K$ with $S(\beta_{k+1}) = \beta_k$, and let $H$ be their span over $\F_2$. So $H$ has $N = 2^m$ points, and $w_m$ is the polynomial of degree $N$ that vanishes exactly on $H$. In our instance $m = 22$.

\head{Rows and cells.} Put $\theta = \beta_1$. Then $\theta^2 = \theta + 1$ and $\{0, 1, \theta, \theta^2\}$ is the field $\F_4$. Let $T$ be the span of $\beta_2, \dots, \beta_{m-1}$, with $n = N/4$ points and vanishing polynomial $v_T$. Every point of $H$ is $x + c$ for exactly one $x \in T$ and one $c \in \F_4$, so $H$ splits into $n$ \emph{rows} of four \emph{cells}:
\begin{center}
\begin{tabular}{@{}l@{\qquad}cccc@{}}
\toprule
cell of the row of $x \in T$ & $\cX$ & $\cY$ & $\cC$ & $\cV$ \\
\midrule
its point & $x$ & $x + 1$ & $x + \theta$ & $x + \theta^2$ \\
\bottomrule
\end{tabular}
\end{center}
The polynomial $r(X) = X^4 + X$ names the row: $r(x + c) = r(x)$ for every $c \in \F_4$.

\section{What is proved}

A \emph{witness} gives a value $w(x)$ to each of the $N$ cells. The circuit is a permutation $\sigma$ of $H$ and eight weight functions $\lambda_0, \dots, \lambda_7$ on $H$. The statement is eight field elements $p_0, \dots, p_7$. The prover shows it knows a witness with three properties.

\head{1. Every row is a full adder.} An element $u = u_0 + \theta u_1$ of $\F_4$ stores two bits, $\bit(u) = u_0$ and $\car(u) = u_1$, and both are polynomials in $u$:
\[
  \bit(u) = \theta u^2 + \theta^2 u, \qquad \car(u) = u^2 + u .
\]
A row adds the bits of its $\cX$ and $\cY$ cells and the carry of its $\cC$ cell, and writes the sum bit and the carry out in its $\cV$ cell. With $\xi = \bit(w_\cX)$, $\eta = \bit(w_\cY)$, $\zeta = \car(w_\cC)$:
\begin{equation}\label{eq:G}
  G(w_\cX, w_\cY, w_\cC, w_\cV) \;=\; w_\cV + (\xi + \eta + \zeta) + \theta \, (\xi \eta + \zeta \, (\xi + \eta)) \;=\; 0 .
\end{equation}
This is the only gate, and there are no selectors: an XOR or an AND is a row that ignores half of its output, and a 32-bit addition is 32 rows, the carry riding from one $\cV$ cell to the next $\cC$ cell. The equation is imposed on arbitrary field elements: that every cell is in $\F_4$ is not assumed, it follows from the three properties together when the circuit has no loop.

\head{2. Copies.} $w(\sigma(x)) = w(x)$ for every cell $x$. The cycles of $\sigma$ are the groups of cells that must be equal: an output and the cells that read it.

\head{3. Public words.} $\sum_{x \in H} \lambda_J(x) \, w(x) = p_J$ for each $J$. This ties chosen cells to the statement.

\section{Two ideas}

\head{Four values open a whole row.} Let $W$ be the polynomial that takes the value $w(x)$ at each cell. Write it on the basis $X_0 = 1$, $X_1 = X$, $X_2 = X^2 + X$, $X_3 = X^3 + X^2$ with coefficients that are polynomials in the row:
\[
  W(X) = \sum_{j=0}^{3} X_j(X) \; g_j\bigl(r(X)\bigr) .
\]
The witness is these four polynomials $g_0, \dots, g_3$. Since $r$ does not move inside a row, the four numbers $a_j = g_j(r(t))$ give $W$ at all four cells of the row of any point $t$, even outside $H$:
\begin{equation}\label{eq:row}
  W(t + c) = \sum_{j=0}^{3} X_j(t + c) \; a_j , \qquad c \in \F_4 .
\end{equation}
So the gate, which reads four cells, costs four field elements and one opening.

\head{Copies are a sum of fractions.} For random $\alpha, \beta$ let
\[
  D(X) = \alpha + W(X) + \beta X, \qquad D'(X) = \alpha + W(X) + \beta \, \sigma(X),
\]
where $\sigma(X)$ is the polynomial that carries the permutation on $H$. If $w(\sigma(x)) = w(x)$ then $D'(x) = D(\sigma(x))$, so the terms $1/D'$ are the terms $1/D$ in another order, and in characteristic 2
\[
  \sum_{x \in H} \Bigl( \frac{1}{D(x)} + \frac{1}{D'(x)} \Bigr) = 0 .
\]
If some copy fails, this sum is nonzero for almost all $\alpha, \beta$. Adding the public words to the same sum costs nothing: with a third challenge $\gamma$ and $\Lambda_\gamma(x) = \sum_J \gamma^J \lambda_J(x)$,
\begin{equation}\label{eq:target}
  \sum_{x \in H} \Bigl( \frac{1}{D(x)} + \frac{1}{D'(x)} + \Lambda_\gamma(x) \, w(x) \Bigr) = \tau, \qquad \tau = \sum_{J} \gamma^J p_J ,
\end{equation}
holds exactly when the copies and the eight public words do, up to a negligible error.

\section{The protocol}

\begin{center}
\fbox{\begin{minipage}{0.94\textwidth}
\vspace{0.3em}
\renewcommand{\arraystretch}{1.35}
\begin{tabular}{@{\ }l@{\quad}p{0.86\textwidth}@{}}
\textbf{P} & sends $g_0, g_1, g_2, g_3$. \\
\textbf{V} & sends $\alpha, \beta, \gamma \in \E$. \\
\textbf{P} & sends $h$, $q$ and $q_G$. \\
\textbf{V} & sends $t \in \E$. \\
\textbf{P} & sends $a_0, a_1, a_2, a_3$ and $s$. \\
\textbf{V} & sends $\rho \in \E$, then checks four evaluation claims.
\end{tabular}
\vspace{0.3em}
\end{minipage}}
\end{center}

\head{The second message.} $h$ is the summand of \eqref{eq:target} as a polynomial of degree below $N$, and $q$, $q_G$ are two quotients. The honest prover's satisfy
\begin{align}
  h \cdot D D' + q \cdot w_m &= D + D' + \Lambda_\gamma \, W \, D D' , \label{eq:copy} \\[2pt]
  q_G \cdot v_T &= G\bigl(W(X), \, W(X + 1), \, W(X + \theta), \, W(X + \theta^2)\bigr) , \label{eq:gate} \\[2pt]
  \textstyle\sum_{x \in H} h(x) &= \tau . \label{eq:sum}
\end{align}
The first says $h = 1/D + 1/D' + \Lambda_\gamma W$ on $H$. The second says the gate holds on every row. The third is \eqref{eq:target}.

\head{The third message.} $a_j = g_j(r(t))$ and $s = \sigma(t)$.

\head{What the verifier computes.} From $t$, the $a_j$ and $s$ alone: the four cells of the row of $t$ by \eqref{eq:row}, hence $W(t)$ and the gate value $G(t)$; then $d = \alpha + W(t) + \beta t$ and $d' = \alpha + W(t) + \beta s$; and the public quantities $\Lambda_\gamma(t)$, $v_T(t)$, $w_m(t)$, $\tau$. Let
\[
  \Delta = d \, d' + w_m(t), \qquad e = d + d' + \Lambda_\gamma(t) \, W(t) \, d \, d' .
\]
It rejects if $\Delta \cdot v_T(t) = 0$.

\head{The four claims.}
\begin{center}
\renewcommand{\arraystretch}{1.5}
\begin{tabular}{@{}cll@{}}
\toprule
 & about & claim \\
\midrule
1 & $g_0, \dots, g_3$ & $\displaystyle \sum_j \kappa_j \, g_j(r(t)) = \sum_j \kappa_j \, a_j$, \quad $\kappa = \bigl((1 + \rho)(1 + \rho^2), \; \rho (1 + \rho^2), \; (1 + \rho) \rho^2, \; \rho^3\bigr)$ \\
2 & $\sigma$ & $\sigma(t) = s$ \\
3 & $h, q, q_G$ & $\displaystyle (1 + \rho) \, \frac{d \, d' \, h(t) + w_m(t) \, q(t)}{\Delta} + \rho \, \frac{d \, d' \, q_G(t)}{\Delta} = \frac{(1 + \rho) \, e \, v_T(t) + \rho \, d \, d' \, G(t)}{\Delta \, v_T(t)}$ \\
4 & $h$ & $\displaystyle \sum_{x \in H} h(x) = \tau$ \\
\bottomrule
\end{tabular}
\end{center}
Each claim is one linear combination of the committed polynomials at one point, equal to a number the verifier has computed. Claim 3 is \eqref{eq:copy} and \eqref{eq:gate} at $t$, added with weights $1 + \rho$ and $\rho$ and divided by $\Delta \, v_T(t)$: the values $h(t)$, $q(t)$, $q_G(t)$ are never sent. Claim 4 is also a single coefficient of $h$ (Section~\ref{sec:pcs}).

\section{Why it is sound}

Suppose the four claims hold of the polynomials the prover committed to.

\begin{itemize}
  \item Claim 1 is an equality of two polynomials of degree 3 in $\rho$, chosen after the $a_j$. So $a_j = g_j(r(t))$, and by \eqref{eq:row} the verifier's cells are $W$ on the row of $t$. Claim 2 gives $s = \sigma(t)$.
  \item Claim 3 is then \eqref{eq:copy} and \eqref{eq:gate} at $t$, separated by $\rho$. Both are polynomial identities of degree $d_{\max} \le 2^{27}$ checked at a random point, so they hold as polynomials.
  \item \eqref{eq:gate} makes every row a full adder. \eqref{eq:copy} with claim 4 gives \eqref{eq:target}, hence the copies and the public words.
\end{itemize}
The three challenges fail with probabilities at most
\[
  \frac{2N^2 + 4N + 8}{|\E|} \ \ (\alpha, \beta, \gamma), \qquad \frac{d_{\max}}{|\E|} \ \ (t), \qquad \frac{3}{|\E|} \ \ (\rho),
\]
each below $2^{-128}$. The polynomials are bounded in degree by the size of what is committed: $g_j$ below $2n$, $h$ below $N$, $q$ and $q_G$ below $4N$.

\section{Zero knowledge}

The prover adds to each $g_j$ a random multiple of the polynomial vanishing on the rows:
\[
  g_j(X) + w_{m-2}(X) \cdot \bigl(\mu_0 + \mu_1 X + \mu_2 (X^2 + X)\bigr), \qquad \mu_0, \mu_1, \mu_2 \in \K \text{ uniform} .
\]
This changes no cell. At a challenge $t$ whose row $r(t)$ is outside $\K$ and off the rows, the mask's value is a uniform element of $\E$, because $1, r, r^2$ is a basis of $\E$ over $\K$. So the four $a_j$ are uniform, and $s$ is public. The verifier sees nothing else of the polynomials except that the four claims are true, which is for the commitment to hide.

\section{With a multilinear commitment}\label{sec:pcs}

Call the \emph{table} of a polynomial $p$ of degree below $2^k$ its coefficients $(c_i)$ on the products of the $w$'s, $p = \sum_i c_i \prod_j w_j^{\,i_j}$ where $i_j$ are the bits of $i$. Then for any $t$,
\[
  \widetilde{c}\,(u_0, \dots, u_{k-1}) = \frac{p(t)}{\prod_{j < k} (1 + w_j(t))} \qquad \text{at } u_j = \frac{w_j(t)}{1 + w_j(t)} ,
\]
where $\widetilde{c}$ is the multilinear extension of the table. So an evaluation of $p$ is one multilinear claim on its table, and the sum $\sum_{x \in H} p(x)$ is the last entry $c_{2^k - 1}$, the claim at the point $(1, \dots, 1)$. Extra coordinates stack several polynomials in one table and combine them linearly, which is how each of the four claims above is a single opening. Any multilinear commitment over these fields can be used, and a table known to vanish outside fixed chunks can be committed as those chunks alone, its claims unchanged.

\end{document}
